\pagebreak
\section{Group}

We introduced the notion of groups in Section 3.1 and we will build on this idea and explore further geometrical applications.

\defn{Order of Group}{
    For a gorup \(G\) with a finite number of elements, the order of the group is the number of elements in \(G\). We write this as \(|G|\).
}

\begin{theorem}
    Let \(G\) be a group. Then
    \begin{enumerate}[label=\alph*)]
        \item The identity element is unique.
        \item For any \(g \in G\), the inverse \(g^{-1}\) is unique.
        \item We can cancel elements if \(ab = ac\) then \(b = c\).
    \end{enumerate}
\end{theorem}

\subsection{Small Finite Groups}

\defn{Cyclic Group}{
    If \(n\) is a positive integer we can define an abstract group of order \(n\) called the cyclic group, \(C_n\), by
    \[C_n = \langle a \mid a^n = e\rangle = \{e, a, a^2, \dots, a^{n - 1}\},\]
    with \(a^n = e\).
}

\begin{proposition}
    All cyclic groups are abelian.
\end{proposition}

\subsection{Group Morphisms}

\defn{Isomorphism}{
    A map \(\phi: G \to H\) between two groups is an isomorphism if it is a bijection such that for all \(g_1, g_2 \in G\),
    \[\phi(g_1 \cdot g_2) = \phi(g_1) * \phi(g_2).\]
    The groups \(G\) and \(H\) are then isomorphic.
}

\begin{proposition}
    If \(G\) and \(H\) are finite and isomorphic, then \(|G| = |H|\).
\end{proposition}

\defn{Homomorphism}{
    A map \(\phi: G \to H\) between two groups is a homomorphism if for all \(g_1, g_2 \in G\)
    \[\phi(g_1 \cdot g_2) = \phi(g_1) * \phi(g_2).\]
}

\begin{theorem}
    Let \(G\) and \(H\) be groups and \(\phi: G \to H\) a homomorphism. Then
    \begin{enumerate}[label=\alph*)]
        \item \(\phi(e_G) = e_H\)
        \item For all \(g \in G, \phi(g^{-1}) = (\phi(g))^{-1}\).
        \item For all \(g \in G, \phi(g^n) = (\phi(g))^n\) for all \(n \in \bZ\).
        \item The image of \(\phi, \operatorname*{im}(\phi)\), is a subgroup of \(H\), and \(\operatorname*{im}(\phi)\) is isomorphic to \(G\) if \(\phi\) is one-to-one.
    \end{enumerate}
\end{theorem}

\subsection{The Symmetric and Dihedral Groups}

\defn{Permutation}{
    A permutation on a finite set, \(S\), of \(n\) distinct objects is a bijection \(\sigma: S \to S\).
}

\begin{theorem}
    The set of all permutations on \(\{1, 2, 3, \dots, n\}\) is a group of order \(n!\).
\end{theorem}

\defn{Symmetric Group}{
    We call the group of all permtuations on \(n\) objects the symmetric group (on \(n\) letters), denoted \(\mathcal{S}_n\).
}

\defn{Dihedral Group}{
    The symmetry group of the regular \(n\)-gon is the group of \(2n\) elements \(D_{2n}\), known as the dihedral group.
}

\begin{theorem}
    Suppose \(G\) is a finite subgroup of linear isometries of the plane. Then \(G\) is isomorphic to \(C_n\) or \(C_{2n}\) for some positive integer \(n\).
\end{theorem}

\subsection{Order and Subgroups}

\defn{Order of Element}{
    For a group \(G\), and an element \(g \in G\), the order of \(g, \mathrm{ord}(g)\), is the smallest positive integer \(n\) for which \(g^n = e\), if such an integer exists.
}

\begin{proposition}
    For any group \(G\) and \(g \in G\) of order \(n\), the set \(H = \{e, g, g^2, \dots, g^{n - 1}\}\) is a (cyclic) subgroup of \(G\) of order \(n\).
\end{proposition}

\defn{Generators}{
    We say that \(H\) in the previous result is generated by \(g\), and usually write the group generated by all powers of an element as \(\langle g \rangle\).
}

\thrm{Lagrange's Theorem}{
    If \(G\) is a finite group and \(H\) is a subgroup of \(G\), then \(|H|\) is a factor of \(|G|\).
}

\defn{Cosets}{
    Let \(G\) be a group and \(H\) a subgroup of \(G\). For any \(g \in G\) we define the left coset of \(H\) by \(g\) to be
    \[gH = \{gh \mid h \in H\}.\]
    There is a corresponding definition of a right coset of \(H\) by \(g\) to be
    \[Hg = \{hg \mid h \in H\}.\]
}

\begin{theorem}
    Let \(g\) be an element of a group \(G\). Then
    \begin{enumerate}[label=\alph*)]
        \item the order of \(g\) is equal to the order of the subgroup \(\langle g \rangle\);
        \item the order of \(g\) is a factor of the order of the group \(|G|\);
    \end{enumerate}
\end{theorem}

\begin{theorem}
    The only group of prime order are the cyclic groups, \(C_p\), and their only subgroups are the trivial ones, \(\{e\}\) and \(C_p\).
\end{theorem}

\begin{corollary}
    All groups of prime order are abelian.
\end{corollary}

\subsection{Normal subgroups and Quotient Groups}

\defn{Normal}{
    A subgroup \(H\) of group \(G\) is normal if for all \(g \in G, gH = Hg\), or equivalently, \(gHg^{-1} = H\).
}

\begin{lemma}
    Let \(H\) be a subgroup of \(G\). If for every \(g \in G\) and \(h \in H\) the element \(g.h.g^{-1}\) is in \(H\), then \(H\) is normal.
\end{lemma}

\begin{proposition}
    Suppose \(G\) is a group of order \(2k\). Then any subgroup of order \(k\) is normal.
\end{proposition}

\begin{theorem}
    Let \(G\) be a subgroup of the isometries \(\mathcal{I_2}\) and \(N\) the subset of translations in \(G\). Then \(N\) is a normal subgroup of \(G\).
\end{theorem}

\begin{theorem}
    Let \(G\) be any group and \(N\) a normal subgroup. Then the set of all cosets, \(\{xN\}\), is a group under the operation
    \[(xN)(yN) = (xy)N.\]
    If \(|G|\) is finite the order of this group is \(|G| / |N|\).
\end{theorem}

\defn{Quotient Groups}{
    The group of cosets in the above theorem is called the factor group or quotient group, and denoted \(G / N\).
}

\begin{theorem}
    If \(N\) is a normal subgroup of \(G\) then the map \(\phi: G \to G / N\) defined by \(\phi(g) = gN\) is a homomorphism.
\end{theorem}

\defn{Natural Homomorphism}{
    The map in the above theorem is called the natural homomorphism from \(G\) to \(G / N\).
}

\defn{Kernel}{
    Let \(\phi: G \to H\) be a homomorphism of groups. The kernel of \(\phi, \ker\phi\), is the set of all elements that maps to the identity (in \(H\)):
    \[\ker(\phi) = \{g \in H \mid \phi(g) = e_H\}.\]
}

\begin{proposition}
    A homomorphism \(\phi: G \to H\) is one-to-one (injective) if and only if \(\ker(\phi) = \{e_G\}\).
\end{proposition}

\begin{theorem}
    The kernel of any homomorphism \(\phi: G \to H\) is a normal subgroup of \(G\).
\end{theorem}

\thrm{First Isomorphism Theorem}{
    For any homomorphism \(\phi: G \to H\), the factor group \(G / \ker(\phi)\) is isomorphic to \(\mathrm{im}(\phi)\).
}

\subsubsection{Group Theory vs Linear Algebra}

Recall that two vector spaces over the same field are isomorphic (as vector spaces) if and only if they have the same dimension. Thus if we are translating group theory to linear algebra, we can replace isomorphic with have the same dimension. \\

In fact, the first isomorphism theorem is just the group theory version of linear algebra's rank-nullity theorem.

\subsection{Product Groups}

\defn{Product Groups}{
    Let \((G, \cdot)\) and \((H, *)\) be two groups. The (direct) product group, is the set of ordered pairs
    \[G \times H = \{(g, h) \mid g \in G, h \in H\}\]
    with group operation
    \[(g_1, h_1)(g_2, h_2) = (g_1 \cdot g_2, h_1 * h_2).\]
}

If \(G\) and \(H\) are finite then clearly \(|G \times H| = |G| \cdot |H|\).

\begin{proposition}
    Let \(G\) and \(H\) be two groups. Then
    \begin{enumerate}[label=\alph*)]
        \item \(G \times H\) and \(H \times G\) are isomorphic.
        \item If \(G\) and \(H\) are abelian, so is \(G \times H\).
    \end{enumerate}
\end{proposition}

\subsubsection{Groups of Order 8}
So far we have
\begin{enumerate}
    \item \(C_8\) --- cyclic (and hence abelian)
    \item \(D_8\) --- the symmetry group of the square; not abelian.
    \item \(C_2 \times C_2 \times C_2\) --- abelian.
\end{enumerate}

One more obvious group of order \(8\) is \(C_4 \times C_2\), also abelian. There is in fact a fifth group of order 8, \(Q_8\), the quaternion group. Its most basic representation is the linear one on \(\bC^2\).
\[Q_8 = \left\{\pm \begin{pmatrix}
        1 & 0 \\ 0 & 1
    \end{pmatrix}, \pm \begin{pmatrix}
        i & 0 \\ 0 & -i
    \end{pmatrix}, \pm \begin{pmatrix}
        0 & -1 \\ 1 & 0
    \end{pmatrix}, \pm \begin{pmatrix}
        0 & i \\ i & 0
    \end{pmatrix}\right\}
\]

But we can more compactly represent the group on the quaternions.
\defn{Quaternions}{
    The quaternions, \(\mathbb{H}\), are the entities formed by taking real linear combinations of the four basis elements \(\{1, i, j, k\}\) where \(i^2 = j^2 = k^2 = ijk = -1\).
}

A quaternion looks like \(a = a_0 + a_i i + a_j j + a_k k\) for real \(a_0, a_i, a_j, a_k\). Multiplication of general quaternions is defined in the obvious way, but it is not commutative. \\

We can now represent \(Q_8\) as the group pf unit quaternions,
\[\{\pm 1, \pm i, \pm j, \pm k\}\]
under multiplication in \(\mathbb{H}\).

\subsubsection{Symmetrices of the Cube}

The physical symmetry group of the cube has order 24 (it would be 48 if we allowed reflections).
\begin{enumerate}[label=\alph*)]
    \item One identity element, of order 1.
    \item 9 face-centred rotations: i.e. through \(\pm \frac{1}{2}\pi\) or \(\pi\) about lines through the centres of the three paris of opposite faces.
    \item 8 vertex-centred rotations: i.e. through \(\pm 2\pi / 3\) around the four space diagonals.
    \item 6 edge-centred rotations: i.e. through \(\pi\) about lines joining midpoints of the six paris of opposite edges.
\end{enumerate}

This group turns out to be isomorphic to \(\mathcal{S}_4\): the group of permutation of 4 objects.

% \subsubsection{Enumeration of Small Groups}

\subsection{Frieze Groups}

We have so far met two principal types of symmetry groups
\begin{enumerate}[label=\alph*)]
    \item finite ones, all realizable as the symmetry group of some geometric object;
    \item infinite ones depending on some continuous parameter, such as rotations about the origin. The symmetry group of a circle is an obvious example.
\end{enumerate}

The finite symmetry groups allow rotations through a discrete set of angles only (e.g. \(\frac{1}{2}\pi\) for a square). \\

There are in fact, patterns whose symmetry groups are infinite but consist of discrete isometries. The simplest sort are the \textbf{frieze groups}, where there is a 2-dimensional pattern repeating on one dimension (a \textbf{frieze}). Such groups have translations in one direction only, and to be discrete, the translations must all be integer multiples of some minimal translation.

\subsection{Wallpaper Group}

Wallpaper groups are the equivalent to frieze groups btu with two independent minimal translations.

\defn{Orbit and Stabliziers}{
    Let \(G\) be a group of transformations in \(\bR^2\) and \(P \in \bR^2\) be any point. \\

    The \textbf{orbit} of \(P\) under \(G, \mathcal{O}(P)\), is the set of all points of \(\bR^2\) you get from applying an arbitrary member of \(G\) to \(P\). The \textbf{stabilizer} of \(P, G_P\), is the set of elements of \(G\) that fix \(P\).
}

Clearly for any two points \(X\) and \(Y\) in the orbit of \(P\) there is an element of \(G\) taking \(X\) to \(Y\). It is also not hard to see that the stabilizer of a point is a subgroup of \(G\) and the same (isomorphic) at every point of a given orbit. It is also called the \textbf{isotropy group} of a point (or orbit).

\begin{theorem}
    For a wallpaper group \(G\) with translation subgroup \(N, G / N\) is isomorphic to \(G_O\), the stabilizer of the origin.
\end{theorem}

\begin{theorem}
    For a wallpaper group \(G\), the isotropy group of the origin, \(G_O\) is \(C_n\) or \(D_{2n}\) for some integer \(n\).
\end{theorem}

\thrm{The Crystallographic Restriction Theorem}{
    Let the symmetry group, \(G\), of a 2-dimensionallattice admit both translations and rotations about one of its lattic points. Then those rotations can only be of period \(1, 2, 3, 4\) or \
    6 i.e. through multiples of \(2\pi / n\) for \(n = 1, 2, 3, 4\) or \(6\).
}

\subsection{Tilings}

Closely related to the idea of wallpaper groups is the concept of a tessellation or tiling of the plane: covering the plane completely with copies of one (or more) geometric shapes (``tiles'').

\begin{lemma}
    The composition of two half turns around point \(P\) and then \(Q\) is a translation through \(2\overrightarrow{PQ}\).
\end{lemma}